% ============================================================
%  Section X.X  Longitudinal User Trajectory Model
%  Paste into the Methods chapter after Section 6.2 (FE Panel)
% ============================================================

\subsection{Longitudinal User Trajectory Model}
\label{sec:longitudinal}

The fixed-effects panel regression of Section~\ref{sec:fe} captures
within-user variation by asking whether a user's lexical quality is
higher in months when they post more frequently.  It does not, however,
address a different and arguably more fundamental question: does
sustained engagement with Reddit reshape a user's language over the
long run?  To test this, we propose a longitudinal growth-curve model
that replaces the calendar time axis with a \emph{user-specific
exposure clock}, tracking each user from their first observed post
and modeling lexical quality as a function of their cumulative
activity history.

\paragraph{Cumulative exposure variables.}
For user $u$ observed in subreddit $s$ during month $t$, define the
two exposure measures

\begin{equation}
  C_{u,t} \;=\; \sum_{\tau \,\leq\, t} n_{u,\tau},
  \qquad
  T_{u,t} \;=\; t - t_{u}^{(0)},
  \label{eq:exposure}
\end{equation}

where $n_{u,\tau}$ is the number of comments posted by $u$ in month
$\tau$ (equivalently, \texttt{num\_utterances\_by\_speaker\_month}
in the combined dataset), and $t_{u}^{(0)}$ is the index of the
earliest month in which $u$ appears.  $C_{u,t}$ therefore counts
the total volume of Reddit activity accumulated by $u$ through month
$t$, while $T_{u,t}$ measures the elapsed \emph{tenure} on the
platform, irrespective of posting volume.  Including both allows us
to disentangle the effect of total output from the mere passage of
time.  Following the convention established in Section~\ref{sec:fe},
cumulative volume enters the model in log-scale as $\log(1 +
C_{u,t})$ to handle the heavy right-tail of the post-count
distribution.

\paragraph{Model specification.}
Retaining the notation of Sections~\ref{sec:wls}--\ref{sec:me}, the
longitudinal trajectory model is

\begin{equation}
  y_{u,s,t}
  \;=\;
  \underbrace{\bigl(\mu_{\alpha} + \alpha_{u}\bigr)}_{\text{user intercept}}
  \;+\;
  \underbrace{\bigl(\mu_{\beta} + b_{u}\bigr)}_{\text{user slope}}
  \log\!\bigl(1 + C_{u,t}\bigr)
  \;+\;
  \gamma\,T_{u,t}
  \;+\;
  \delta\,x_{u,s,t}
  \;+\;
  \boldsymbol{\eta}^{\!\top}\mathbf{z}_{u,s,t}
  \;+\;
  \varepsilon_{u,s,t},
  \label{eq:longitudinal}
\end{equation}

where

\begin{itemize}
  \item $\mu_{\alpha}$ is the population-level intercept and
        $\alpha_{u} \sim \mathcal{N}(0,\,\sigma_{\alpha}^{2})$ is a
        user-level random intercept absorbing all time-invariant
        heterogeneity (e.g.\ baseline vocabulary, educational
        background, native language);

  \item $\mu_{\beta}$ is the \emph{average} effect of log-cumulative
        exposure on lexical quality---the parameter of primary
        interest---and $b_{u} \sim \mathcal{N}(0,\,\sigma_{b}^{2})$
        is a user-level random slope that allows individual
        trajectories to diverge from the population mean;

  \item $[\alpha_{u},\, b_{u}]^{\top} \sim
        \mathcal{N}\!\left(\mathbf{0},\,\boldsymbol{\Sigma}\right)$
        with $\boldsymbol{\Sigma}$ estimated freely, permitting
        correlation between a user's baseline quality and their
        rate of change;

  \item $\gamma$ captures any residual effect of tenure net of
        cumulative posting volume;

  \item $x_{u,s,t} = \log(1 + n_{u,t})$ is the current-month
        log-comment count (as defined in Section~\ref{sec:wls}),
        included to control for contemporaneous posting intensity;

  \item $\mathbf{z}_{u,s,t}$ is the vector of comment-level
        controls---post depth, edited status, Reddit score, number
        of direct replies, and controversiality---defined in
        Section~\ref{sec:wls};

  \item $\varepsilon_{u,s,t} \overset{\mathrm{iid}}{\sim}
        \mathcal{N}(0,\,\sigma^{2})$.
\end{itemize}

The model is estimated separately for each of the six lexical quality
metrics via restricted maximum likelihood (REML), fitting one model
per metric pooled across all subreddits with subreddit fixed effects
absorbed into the intercept.  Only users with at least
$m_{\min} = 6$ months of observed data are included, ensuring that
individual trajectories are identified.  $p$-values for the key
coefficients $\mu_{\beta}$ and $\gamma$ are adjusted using the
Benjamini--Hochberg procedure across the six simultaneous tests.

\paragraph{Interpretation.}
A significantly positive $\hat{\mu}_{\beta}$ for a given metric
indicates that, on average, users' lexical quality on that dimension
grows with cumulative Reddit exposure, consistent with a
practice-makes-perfect hypothesis.  A significantly negative value
would instead suggest that heavier long-run use is associated with
declining quality.  The variance component $\sigma_{b}^{2}$ speaks
to how heterogeneous these trajectories are: large $\sigma_{b}^{2}$
relative to $\hat{\mu}_{\beta}$ implies that some users improve
while others decline, and that aggregate conclusions should be
interpreted cautiously.  The coefficient $\gamma$ isolates any
effect of time-on-platform beyond posting volume---for instance,
passive acculturation to community norms---while $\delta$ controls
for the short-run activity effect already characterized in
Section~\ref{sec:fe}.
